\chapter{Error Handling}
\index{error handling}\index{exceptions}
\label{ch:errors}

Programs encounter errors: files that do not exist, invalid user input, network failures, division by zero. Scheme provides a structured exception system with \texttt{guard} for catching errors, \texttt{raise} for signaling them, and \texttt{error} for creating error objects. This chapter covers all three, plus strategies for writing robust code. When you hit an error you do not recognize, Appendix~\ref{app:errors} lists the most common messages along with their causes and fixes.

\section{The \texttt{error} Procedure}
\index{error}

The simplest way to signal an error is the \texttt{error} procedure:

\begin{lstlisting}[style=repl]
> (error "something went wrong")
error[KP3000]: something went wrong
\end{lstlisting}

\texttt{error} creates an \emph{error object}\index{error object} and raises it as an exception. The first argument is a message string; additional arguments are \emph{irritants}\index{irritants}---values that provide context about what went wrong, as in \texttt{(error "index out of range" 5)}.

At the top level, the REPL prints the message and irritants and abandons the evaluation. They are not just display text: they travel inside the error object, and any handler that catches it can read them back---as \texttt{guard} does in the next section. The \texttt{[KP3000]} tag is a \emph{diagnostic code}: every error Kaappi reports carries one, and Section~\ref{sec:error-codes} shows how to put them to work.

\subsection{Using \texttt{error} in Your Functions}

\begin{lstlisting}[style=repl]
> (define (safe-divide a b)
    (if (zero? b)
        (error "division by zero" a b)
        (/ a b)))
> (safe-divide 10 2)
5
> (safe-divide 10 0)
error[KP3000]: division by zero 10 0
\end{lstlisting}

Note how the irritants \texttt{10} and \texttt{0} appear after the message---this is why you pass the offending values to \texttt{error} instead of formatting them into the string.

\subsection{What an Uncaught Error Looks Like}

When no handler catches an error, Kaappi reports it and abandons the evaluation. Errors detected by the runtime itself---type errors, arity mismatches, undefined variables---are reported with their full detail, and running a script also gets you the file, line, and column of the offending expression, the source line itself, and the call site it was reached from:

\begin{Verbatim}[fontsize=\footnotesize]
$ kaappi process.scm
process.scm:2:3: error[KP3002]: type error in 'car': expected pair, got 5
      (car n))
  called from process.scm:3
\end{Verbatim}

Learn to read these reports before learning to catch errors---most of the time, the right response to an unexpected error is to fix the bug it points at, not to handle it.

\begin{note}[Coming from Python]
\texttt{(error "message" value)} is like Python's \texttt{raise ValueError(f"message: \{value\}")}. The key difference is that Scheme's error objects are structured---the message and irritants are separate fields, not formatted into a string.
\end{note}

\section{Catching Errors with \texttt{guard}}
\index{guard}

\texttt{guard} is Scheme's primary exception-handling form. It catches raised exceptions and matches them against clauses, similar to \texttt{try}/\texttt{except} in Python:

\begin{lstlisting}[style=repl]
> (guard (e (#t (string-append "caught: "
                               (error-object-message e))))
    (error "file not found" "data.txt"))
"caught: file not found"
\end{lstlisting}

The \texttt{\#t} clause here matches \emph{any} raised object---fine for a first demonstration, but deliberately blunt. Real handlers should name the errors they expect, as the rest of this chapter shows (and as the warning on catch-all handlers below insists).

The general form:

\begin{lstlisting}
(guard (variable
        (test1 expr1 ...)
        (test2 expr2 ...)
        ...
        (else expr ...))
  body ...)
\end{lstlisting}

The \texttt{body} is evaluated. If an exception is raised, it is bound to \texttt{variable}, and the clauses are tested in order (like \texttt{cond}). The first matching clause's expressions are evaluated and returned. If no clause matches, the exception is re-raised.

\subsection{Multiple Clauses}

\begin{lstlisting}[style=repl]
> (guard (e
          ((string? e)
           (string-append "string error: " e))
          ((number? e)
           (* e 10))
          (else "unknown error"))
    (raise "oops"))
"string error: oops"
\end{lstlisting}

\subsection{Inspecting Error Objects}

Error objects created by \texttt{error} can be examined with three procedures:

\begin{lstlisting}[style=repl]
> (guard (e (#t (list
                  (error-object? e)
                  (error-object-message e)
                  (error-object-irritants e))))
    (error "bad value" 42 'expected-string))
(#t "bad value" (42 expected-string))
\end{lstlisting}

\begin{itemize}
    \item \texttt{error-object?}\index{error-object?} --- tests if the value is an error object
    \item \texttt{error-object-message}\index{error-object-message} --- extracts the message string
    \item \texttt{error-object-irritants}\index{error-object-irritants} --- extracts the list of irritant values
\end{itemize}

\subsection{\texttt{file-error?} and \texttt{read-error?}}
\index{file-error?}\index{read-error?}

R7RS singles out two error classes with dedicated predicates. \texttt{file-error?} recognizes errors raised by file operations---most commonly, opening a file that does not exist. \texttt{read-error?} recognizes errors raised by \texttt{read} on malformed input:

\begin{lstlisting}[style=repl]
> (guard (e ((file-error? e) 'no-such-file))
    (open-input-file "missing.txt"))
no-such-file
> (guard (e ((read-error? e) 'bad-syntax))
    (read (open-input-string ")")))
bad-syntax
\end{lstlisting}

Prefer these predicates over matching on message text. Error messages differ between Scheme implementations---and may change between Kaappi versions---but the predicates are part of the standard.

\subsection{Conditional Error Handling}

You can match on the error message or type to handle different errors differently:

\begin{lstlisting}
(define (safe-lookup key table)
  (guard (e
          ((and (error-object? e)
                (string=? (error-object-message e)
                          "key not found"))
           #f)
          (else (raise e)))  ; re-raise unknown errors
    (lookup key table)))
\end{lstlisting}

The \texttt{(else (raise e))} clause ensures that unexpected errors are not silently swallowed.

Matching on the message string works here because \texttt{lookup} is our own code and we control the text. It is brittle for anything else: an error message is not part of a procedure's contract and may change without notice. For errors from other code, dispatch on a predicate instead---\texttt{file-error?}, \texttt{read-error?}, a custom error type as shown later in this chapter---or, for errors raised by the interpreter itself, on the stable code shown next.

\begin{warning}[Always re-raise unknown errors]
Do not write \texttt{(guard (e (\#t \#f)) ...)} to blindly catch everything. This hides bugs. Only catch the specific errors you know how to handle, and re-raise the rest.
\end{warning}

\subsection{Stable Error Codes}
\label{sec:error-codes}
\index{error codes}\index{error-object-code}\index{diagnostics library@\texttt{(kaappi diagnostics)}}

Every diagnostic Kaappi reports carries a stable code---the \texttt{KP3002} in \texttt{error[KP3002]: type error in 'car'}---and the built-in \texttt{(kaappi diagnostics)} library reads it back programmatically. \texttt{error-object-code} returns the code as a symbol, or \texttt{\#f} for a value that has none:

\begin{lstlisting}[style=repl]
> (import (kaappi diagnostics))
> (guard (e (#t (error-object-code e)))
    (car '()))
KP3002
> (guard (e (#t (error-object-code e)))
    (error "bad input" 42))
#f
\end{lstlisting}

Note the second result: codes identify \emph{interpreter} diagnostics. Error objects you create with \texttt{error} carry no code---the \texttt{KP3000} the printer shows for them means only ``user-raised error''---so inside your own code you still dispatch on messages or custom error types. Codes shine when you must catch one specific runtime failure without depending on its message text:

\begin{lstlisting}[style=repl]
> (guard (e ((eq? (error-object-code e) 'KP3004)
             'caught-division))
    (/ 1 0))
caught-division
\end{lstlisting}

\texttt{error-object-code} is total---it accepts any value, error object or not, and never raises---so it is safe as the first test in a \texttt{guard} clause. And when a code in a report means nothing to you, ask the interpreter what it stands for:

\begin{Verbatim}[fontsize=\footnotesize]
$ kaappi explain KP3004
KP3004  division-by-zero
runtime · error

division by zero

An exact division, 'modulo', 'remainder', or 'quotient' had a zero
divisor. Guard the divisor, or use inexact arithmetic where the result
is a floating-point infinity or NaN instead of an error.
\end{Verbatim}

Appendix~\ref{app:errors} walks through the most common codes; \texttt{kaappi explain --all} lists every one.

\section{Raising Exceptions with \texttt{raise}}
\index{raise}

\texttt{raise} signals an exception. Any value can be raised---it does not have to be an error object:

\begin{lstlisting}[style=repl]
> (guard (e (#t e))
    (raise "a string error"))
"a string error"
> (guard (e (#t e))
    (raise 42))
42
> (guard (e (#t e))
    (raise '(custom error data)))
(custom error data)
\end{lstlisting}

However, the convention is to use \texttt{error} (which creates a proper error object) for most situations. Direct \texttt{raise} with arbitrary values is useful for custom condition types.

\subsection{\texttt{raise} vs.\ \texttt{error}}

\begin{center}
\begin{tabular}{lp{6cm}}
\textbf{Form} & \textbf{Use when} \\
\hline
\texttt{(error msg irritants ...)} & Signaling an error with a message---the common case \\
\texttt{(raise value)} & Signaling a non-error condition, or raising a custom condition object \\
\end{tabular}
\end{center}

\section{\texttt{with-exception-handler}: Low-Level Handler}
\index{with-exception-handler}

\texttt{with-exception-handler} installs an exception handler for the duration of a thunk:

\begin{lstlisting}[style=repl]
> (call/cc
    (lambda (exit)
      (with-exception-handler
        (lambda (e) (exit (list 'caught e)))
        (lambda () (raise "boom")))))
(caught "boom")
\end{lstlisting}

This example leans on \texttt{call/cc}, which is not covered until Chapter~\ref{ch:continuations}. For now, read \texttt{(call/cc (lambda (exit) ...))} as ``capture an early-return procedure named \texttt{exit}''---calling \texttt{exit} abandons the rest of the computation and returns its argument.

\texttt{with-exception-handler} is lower-level than \texttt{guard}. The handler receives the exception and must decide what to do: it can invoke an escape continuation, re-raise, or (for continuable exceptions) return a value.

In most code, \texttt{guard} is preferred because it is clearer and handles re-raising automatically. Use \texttt{with-exception-handler} only when you need the full power (e.g., logging all exceptions without catching them).

\section{Continuable Exceptions}
\index{raise-continuable}

\texttt{raise-continuable} raises an exception that the handler can ``fix'' by returning a value:

\begin{lstlisting}[style=repl]
> (with-exception-handler
    (lambda (e) (* e 10))    ; handler returns a replacement
    (lambda () (+ 1 (raise-continuable 5))))
51
\end{lstlisting}

The handler receives 5, returns 50, and execution continues as if \texttt{raise-continuable} had returned 50. This is an advanced feature---use it for situations where a handler can supply a default or corrected value.

\section{Common Patterns}

\subsection{Default Values on Error}

Here we intentionally catch all errors---any failure (missing file, parse error, or otherwise) means we fall back to the default:

\begin{lstlisting}
(define (read-config-or-default filename)
  (guard (e (#t '()))  ; return empty list on any error
    (call-with-port (open-input-file filename)
      (lambda (port) (read port)))))
\end{lstlisting}

\subsection{Guaranteed Cleanup}

For guaranteed resource cleanup (similar to Python's \texttt{try}/\texttt{finally}), Scheme provides \texttt{dynamic-wind}\index{dynamic-wind}. Its third thunk runs no matter how the body exits---normally or via an exception:

\begin{lstlisting}[style=repl]
> (guard (e (#t (display "caught\n")))
    (dynamic-wind
      (lambda () (display "enter\n"))
      (lambda () (error "boom"))
      (lambda () (display "cleanup\n"))))
enter
cleanup
caught
\end{lstlisting}

Chapter~\ref{ch:continuations} explains \texttt{dynamic-wind} in full, including why it takes three thunks. For the common case of file handling, the simplest approach remains \texttt{call-with-port}, which closes the port when the body returns or raises an exception.

\subsection{Wrapping External Calls}

\begin{lstlisting}
(define (safe-json-parse str)
  (guard (e
          ((error-object? e)
           (values #f (error-object-message e)))
          (else
           (values #f "unknown parse error")))
    (values (json-read (open-input-string str)) #f)))
\end{lstlisting}

This returns two values: the parsed result (or \texttt{\#f} on error) and an error message (or \texttt{\#f} on success).

\subsection{Asserting Preconditions}

\begin{lstlisting}
(import (srfi 132))   ; provides list-sort

(define (vector-median v)
  (when (zero? (vector-length v))
    (error "cannot compute median of empty vector" v))
  (let ((sorted (list-sort < (vector->list v))))
    (list-ref sorted (quotient (length sorted) 2))))
\end{lstlisting}

Fail early with a clear message when preconditions are not met.

\section{Custom Error Types}
\index{custom error types}

For larger programs, you may want to distinguish between different kinds of errors. One approach is to use records as error objects and raise them with \texttt{raise}. We use \texttt{define-record-type} (covered in Chapter~\ref{ch:records}) to define a custom error type with a type name, a constructor, a predicate, and field specifications:

\begin{lstlisting}
(define-record-type <http-error>
  (make-http-error status message)
  http-error?
  (status http-error-status)
  (message http-error-message))

(define (fetch url)
  (let ((response (http-get url)))
    (unless (= (response-status response) 200)
      (raise (make-http-error
               (response-status response)
               (response-body response))))
    (response-body response)))
\end{lstlisting}

Catch them with \texttt{guard} using the record predicate:

\begin{lstlisting}[style=repl]
> (guard (e
          ((http-error? e)
           (string-append "HTTP "
             (number->string (http-error-status e))
             ": " (http-error-message e)))
          ((error-object? e)
           (error-object-message e)))
    (fetch "https://example.com/missing"))
"HTTP 404: Not Found"
\end{lstlisting}

This pattern gives you type-safe error dispatch---the \texttt{guard} clauses test each error type in order, and unrecognized errors fall through to the next clause or propagate up.

\section{Error Handling in Libraries}

When writing a library, follow these conventions so that callers can handle your errors cleanly:

\begin{itemize}
    \item \textbf{Use \texttt{error} with a consistent prefix.} Start every error message with the library or procedure name: \texttt{(error "json-read: unexpected token" ch)}. This makes it immediately clear where the error originated.
    \item \textbf{Include the relevant values as irritants.} The caller needs context to diagnose the problem. Include the input that caused the error, the expected type, and the actual value.
    \item \textbf{Document which errors your procedures raise.} If \texttt{json-read} can raise a parse error, say so in the documentation. Callers should not have to read your source code to know what to catch.
    \item \textbf{Do not catch errors you cannot handle.} If a low-level I/O function fails inside your library, let that error propagate. Only catch errors that your library can meaningfully recover from.
\end{itemize}

\begin{lstlisting}
(define (parse-config filename)
  (guard (e ((file-error? e)
             (error "parse-config: file not found" filename)))
    (call-with-port (open-input-file filename)
      (lambda (port) (json-read port)))))
\end{lstlisting}

This wraps the low-level file error in a domain-specific message while preserving the offending filename as an irritant. Note the use of \texttt{file-error?} rather than matching the message text---the predicate keeps working even if the underlying message changes.

\section{Hunting Bugs with the Debugger}
\index{debugger}\index{breakpoints}

\texttt{guard} and \texttt{raise} manage the errors you expect. For the bug you do not understand yet---a procedure returning the wrong value, with no error in sight---the REPL has an interactive debugger. \texttt{,break} sets a breakpoint on any procedure. The next time it is called, evaluation pauses and drops you into a \texttt{debug>} prompt where you can look around:

\begin{lstlisting}
(define (average lst)
  (/ (apply + lst) (length lst)))

(define (report lst)
  (display (average lst))
  (newline))
\end{lstlisting}

\begin{lstlisting}[style=repl]
> ,break average
Breakpoint set on average
> (report '(3 4 8))
Break at average (<repl>:1)
(*@\textbf{debug>}@*) locals
  lst = (3 4 8)
(*@\textbf{debug>}@*) backtrace
> [2] average
  [1] report
  [0] <lambda>
(*@\textbf{debug>}@*) continue
5
\end{lstlisting}

\texttt{locals} shows the paused procedure's bindings---exactly what \texttt{average} received, with no \texttt{display} calls sprinkled through your code. \texttt{backtrace} shows how execution got here, innermost frame first (the \texttt{>} marks the current frame): the expression typed at the REPL called \texttt{report}, which called \texttt{average}. \texttt{continue} resumes the program as if nothing had happened.

From the same prompt, \texttt{step} moves into the next expression, \texttt{next} steps over it, and \texttt{finish} runs to the end of the current procedure. \texttt{watch var} pauses when a variable changes---useful when some piece of state mutates and you cannot see where from. When you are done, \texttt{,delete all} clears every breakpoint. The full command set for both the REPL and the \texttt{debug>} prompt is tabulated in Appendix~\ref{app:cli}.

The debugger complements, rather than replaces, the habits from this chapter: read the error report first (it usually names the culprit), reproduce the problem with the smallest input you can, and only then reach for a breakpoint to watch the wrong value being computed.

\section{Error Handling Strategy}

\begin{itemize}
    \item \textbf{Fail fast.} Use \texttt{error} to signal problems as soon as they are detected. Do not pass invalid data deeper into the program hoping something will work.
    \item \textbf{Catch at boundaries.} Handle errors at the boundary between subsystems---at the top of a request handler, at the entry point of a CLI tool, at the edge of a library call. Do not pepper every function with \texttt{guard}.
    \item \textbf{Be specific.} Catch only the errors you can handle. Let unexpected errors propagate---they usually indicate bugs that should not be silently swallowed.
    \item \textbf{Provide context.} Use irritants in \texttt{error} to include the values that caused the problem. ``index out of range'' is less useful than ``index out of range: 10, list length: 5.''
    \item \textbf{Log, then re-raise.} When you need to observe errors without handling them, install a handler that logs and re-raises:
\end{itemize}

\begin{lstlisting}
(define (with-error-logging thunk)
  (guard (e (#t (display "ERROR: " (current-error-port))
               (display (if (error-object? e)
                            (error-object-message e)
                            e)
                        (current-error-port))
               (newline (current-error-port))
               (raise e)))
    (thunk)))
\end{lstlisting}

\section*{Summary}

\begin{itemize}
    \item \texttt{error} signals a condition with a message and irritants, while \texttt{raise} signals any object.
    \item \texttt{guard} catches exceptions with
    \texttt{cond}-style clauses, and
    \texttt{with-exception-\allowbreak handler} is the underlying primitive.
    \item Dispatch on predicates, not message text:
    \texttt{file-error?} and \texttt{read-error?} identify
    the standard error classes.
    \item Records make good custom error types, letting handlers distinguish one kind of error from another.
    \item Common patterns include supplying default values on error, guaranteeing cleanup, and asserting preconditions.
    \item For bugs rather than errors, \texttt{,break} pauses a running program at a procedure so you can inspect \texttt{locals} and the \texttt{backtrace} before continuing.
\end{itemize}

\section*{Exercises}

\begin{exercise}[Safe Division]
Write a \code{safe-divide} procedure that takes a numerator, a denominator, and an optional default value. If the denominator is zero, it returns the default (or \code{0} if none is provided) instead of raising an error. Use \code{guard} to catch the error.

\begin{lstlisting}[style=repl]
> (safe-divide 10 3)
10/3
> (safe-divide 10 0)
0
> (safe-divide 10 0 -1)
-1
\end{lstlisting}
\end{exercise}

\begin{exercise}[Retry on Error]
Write a \code{retry} procedure that takes a thunk and a maximum number of attempts. It calls the thunk; if it raises an error, it tries again up to \emph{n} times. If all attempts fail, it re-raises the last error. Print which attempt is being made.

\begin{lstlisting}[style=repl]
> (let ((count 0))
    (retry (lambda ()
             (set! count (+ count 1))
             (if (< count 3)
                 (error "not yet" count)
                 count))
           5))
Attempt 1 failed: not yet
Attempt 2 failed: not yet
3
\end{lstlisting}
\end{exercise}

\begin{exercise}[Validate with Descriptive Errors]
Write a \code{validate-email} procedure that checks a string for basic email validity. Raise descriptive errors with irritants for each failure case: the string must contain exactly one \code{@} character, must have at least one character before the \code{@}, and must have a \code{.} after the \code{@}. Each error message should include the input string as an irritant.

\begin{lstlisting}[style=repl]
> (validate-email "alice@example.com")
#t
> (validate-email "no-at-sign")
error[KP3000]: missing @ in email address "no-at-sign"
> (validate-email "@example.com")
error[KP3000]: empty local part in email address "@example.com"
\end{lstlisting}
\end{exercise}

\begin{exercise}[With-Default]
Write a \code{with-default} procedure that takes a default value and a thunk. It calls the thunk; if any error is raised, it returns the default value instead. This is similar to Python's pattern of \texttt{try: ... except: return default}.

\begin{lstlisting}[style=repl]
> (with-default 0 (lambda () (+ 1 2)))
3
> (with-default 0 (lambda () (error "boom")))
0
> (with-default '() (lambda () (vector-ref #(1 2) 5)))
()
\end{lstlisting}
\end{exercise}

In the next chapter (Chapter~\ref{ch:continuations}), we explore continuations---Scheme's most powerful and unusual control flow mechanism.

\chapterend
